\documentclass[13pt,a4paper]{article}

\usepackage[
  top=0.6in,
  bottom=0.6in,
  left=0.6in,
  right=0.6in
]{geometry}
\usepackage{amsmath,amssymb}
\usepackage{bm}
\usepackage{fancyhdr}

\pagestyle{fancy}
\fancyhf{}
\cfoot{\thepage}
\renewcommand{\headrulewidth}{0pt}
\renewcommand{\footrulewidth}{0pt}

\begin{document}

Bayesian inference concerns learning about unknown quantities from observed data. Bayes' theorem expresses how information from the data updates an initial probability distribution:
\begin{equation}
p(H \mid D)
=
\frac{p(D \mid H)p(H)}{p(D)}.
\end{equation}

Here, $H$ denotes a hypothesis and $D$ denotes observed data. The quantity $p(H)$ is the prior probability of the hypothesis before observing the data. The likelihood $p(D \mid H)$ is the probability of observing the data if the hypothesis is true. The evidence $p(D)$ is the overall probability of the data under all possible hypotheses. It normalises the posterior so that
\[
\int p(H \mid D)\,dH = 1.
\]

The posterior $p(H \mid D)$ is the updated probability distribution after the data have been observed. It combines the prior information with the compatibility of each hypothesis with the data. More generally, if a model is described by a vector of unknown parameters
\[
\bm{\theta}
=
(\theta_1,\theta_2,\ldots,\theta_d)
\]
and the observed data are denoted by $\bm{x}_0$, Bayes' theorem becomes
\begin{equation}
p(\bm{\theta}\mid\bm{x}_0)
=
\frac{
p(\bm{x}_0\mid\bm{\theta})p(\bm{\theta})
}{
p(\bm{x}_0)
}.
\end{equation}

The posterior $p(\bm{\theta}\mid\bm{x}_0)$ describes which parameter values remain plausible after observing $\bm{x}_0$. The prior $p(\bm{\theta})$ describes the parameter distribution before the observation. The likelihood
\[
p(\bm{x}_0\mid\bm{\theta})
\]
measures how compatible the observed data are with a proposed parameter value. The evidence
\[
p(\bm{x}_0)
=
\int p(\bm{x}_0\mid\bm{\theta})p(\bm{\theta})\,d\bm{\theta}
\]
is obtained by averaging the likelihood over the prior. It is the normalising constant for the posterior and can also be interpreted as the probability of obtaining the observed data under the model.

A useful quantity is the likelihood-to-evidence ratio
\begin{equation}
r(\bm{x}_0\mid\bm{\theta})
=
\frac{p(\bm{x}_0\mid\bm{\theta})}{p(\bm{x}_0)}.
\end{equation}
By Bayes' theorem, this is also the posterior-to-prior ratio:
\begin{equation}
r(\bm{x}_0\mid\bm{\theta})
=
\frac{p(\bm{\theta}\mid\bm{x}_0)}{p(\bm{\theta})}.
\end{equation}

The ratio therefore describes how observing $\bm{x}_0$ changes the plausibility of $\bm{\theta}$. If the ratio is greater than one, the observation increases the relative plausibility of that parameter value. If it is less than one, the observation decreases its relative plausibility. Knowing the ratio and the prior is sufficient to reconstruct the posterior:
\begin{equation}
p(\bm{\theta}\mid\bm{x}_0)
=
r(\bm{x}_0\mid\bm{\theta})p(\bm{\theta}).
\end{equation}

In many physical applications, the parameters describe a forward model. Given a parameter value $\bm{\theta}$, the model produces data according to a simulator:
\[
\bm{x}\sim p(\bm{x}\mid\bm{\theta}).
\]
The simulator may represent physical evolution, measurement effects, detector noise, selection effects, or other stochastic processes. It may be straightforward to generate a synthetic dataset, while being difficult or impossible to evaluate the probability density of a particular dataset analytically.

This situation is described as an intractable likelihood. Standard Bayesian methods require the likelihood, either explicitly or through repeated evaluations of it. When the likelihood cannot be evaluated, ordinary posterior-sampling methods cannot be applied directly.

Simulation-based inference (SBI), also called likelihood-free inference, addresses this setting by using simulations instead of explicit likelihood evaluations. Samples are generated by first drawing parameters from the prior and then simulating data:
\begin{equation}
\bm{\theta}^{(i)}\sim p(\bm{\theta}),
\qquad
\bm{x}^{(i)}\sim p(\bm{x}\mid\bm{\theta}^{(i)}).
\end{equation}
The resulting parameter--data pairs,
\[
\left\{
\left(\bm{\theta}^{(i)},\bm{x}^{(i)}\right)
\right\}_{i=1}^{N},
\]
are used to learn information about the posterior or a ratio related to it.

Markov Chain Monte Carlo (MCMC) is a traditional method for sampling from a posterior distribution. It produces a sequence
\[
\bm{\theta}^{(1)},\bm{\theta}^{(2)},\ldots
\]
in which the next parameter value depends on the current value. At each step, a proposal is generated and accepted or rejected according to a rule designed so that the stationary distribution of the chain is the target posterior.

If the likelihood can be evaluated, MCMC can explore the posterior without requiring the posterior normalisation constant. However, each proposed state may require a likelihood evaluation. If likelihood evaluation is expensive or impossible, MCMC becomes impractical. Even when it is usable, MCMC may spend many evaluations in regions of parameter space that have negligible posterior probability.

SBI and MCMC therefore address the same broad task—Bayesian parameter inference—but use different computational strategies. MCMC generates a chain for one observed dataset by repeatedly evaluating a likelihood or an approximation to it. SBI learns an inference function from simulated parameter--data pairs. Once trained, that function can be evaluated rapidly for the observation of interest.

SBI introduces its own difficulties. First, the parameter space may be high-dimensional. If there are $d$ parameters and the posterior occupies a fraction $f$ of the prior range in each direction, then the posterior occupies approximately
\[
f^d
\]
of the prior volume. For example, if the relevant range is one tenth of the prior range in each of $20$ dimensions, the relevant volume is approximately
\[
(0.1)^{20}=10^{-20}
\]
of the original volume. Random simulations from the full prior will then rarely fall in the region that contributes substantially to the posterior.

This is one manifestation of the curse of dimensionality. As the number of dimensions increases, the volume of parameter space grows rapidly, while the available simulations become increasingly sparse. A method attempting to learn the full joint posterior must represent dependence across all parameter dimensions, including parameters that may not be of direct scientific interest.

Second, many physical models contain nuisance parameters. These parameters affect the simulated data and must be accounted for, but they may not be included in the final result. The desired output is often a one- or two-dimensional marginal posterior rather than the complete joint posterior.

Third, simulations may be expensive. A method that wastes most of its simulations in regions incompatible with the observation is inefficient. Finally, the true posterior is usually unavailable, making calibration and validation important. Credible intervals should be tested to determine whether they contain the true parameter values at their stated frequencies.

Neural Ratio Estimation (NRE) approaches SBI by estimating a probability ratio with a classifier rather than directly estimating a probability density. The target ratio may be written as
\begin{equation}
r(\bm{x}\mid\bm{\theta})
=
\frac{p(\bm{x}\mid\bm{\theta})}{p(\bm{x})}.
\end{equation}

The training data consist of two classes of parameter--data pairs. A joint pair is generated by drawing $\bm{\theta}$ from the prior and then generating $\bm{x}$ using that same parameter:
\begin{equation}
(\bm{\theta},\bm{x})
\sim
p(\bm{\theta},\bm{x})
=
p(\bm{\theta})p(\bm{x}\mid\bm{\theta}).
\end{equation}

An independent pair is generated by drawing the parameter and data independently:
\begin{equation}
(\bm{\theta},\bm{x})
\sim
p(\bm{\theta})p(\bm{x}).
\end{equation}

The classifier is trained to distinguish joint pairs from independent pairs. Joint pairs form one class, and independent pairs form the other. The classifier must learn whether the parameter and data are statistically associated or whether they are unrelated.

Let $s(\bm{x},\bm{\theta})$ denote the classifier output, interpreted as the probability that a pair belongs to the joint class. For balanced classes, the optimal classifier satisfies
\begin{equation}
s(\bm{x},\bm{\theta})
=
\frac{
p(\bm{\theta},\bm{x})
}{
p(\bm{\theta},\bm{x})
+
p(\bm{\theta})p(\bm{x})
}.
\end{equation}
The density ratio is then obtained from the classifier odds:
\begin{equation}
\frac{s(\bm{x},\bm{\theta})}
{1-s(\bm{x},\bm{\theta})}
=
\frac{
p(\bm{\theta},\bm{x})
}{
p(\bm{\theta})p(\bm{x})
}
=
\frac{p(\bm{x}\mid\bm{\theta})}{p(\bm{x})}.
\end{equation}

NRE therefore transforms likelihood-free inference into a binary classification problem. A trained classifier estimates how strongly a particular parameter value and dataset are associated.

Marginalisation is required when only a subset of parameters is of interest. Let
\[
\bm{\theta}
=
(\bm{\phi},\bm{\psi}),
\]
where $\bm{\phi}$ denotes parameters of interest and $\bm{\psi}$ denotes nuisance parameters. The marginal posterior for $\bm{\phi}$ is
\begin{equation}
p(\bm{\phi}\mid\bm{x}_0)
=
\int
p(\bm{\phi},\bm{\psi}\mid\bm{x}_0)
\,d\bm{\psi}.
\end{equation}

Marginalisation does not remove nuisance parameters from the model. Their uncertainty remains present through the integration over $\bm{\psi}$. It removes them only from the reported distribution.

Marginal Neural Ratio Estimation (MNRE) trains ratio estimators for selected low-dimensional parameter subsets rather than for the full vector. For example, an estimator can target
\[
r(\bm{x}\mid\phi_i)
\]
for a one-dimensional marginal, or
\[
r(\bm{x}\mid\phi_i,\phi_j)
\]
for a two-dimensional marginal. The nuisance parameters are still sampled when generating the simulations, but they are not included in the classifier input as target parameters. Their effects are therefore incorporated through the distribution of simulated data.

TMNRE extends MNRE by truncating the prior iteratively. The procedure begins with simulations drawn from the original prior. The resulting simulations are used to train marginal ratio estimators. These estimators are then evaluated for the observed data $\bm{x}_0$ to identify regions of parameter space with appreciable posterior probability.

Let $\Gamma$ denote the retained region. The truncated prior is
\begin{equation}
p_{\Gamma}(\bm{\theta})
=
\frac{
p(\bm{\theta})\mathbb{I}_{\Gamma}(\bm{\theta})
}{
\displaystyle
\int_{\Gamma}p(\bm{\theta})\,d\bm{\theta}
},
\end{equation}
where
\[
\mathbb{I}_{\Gamma}(\bm{\theta})
=
\begin{cases}
1, & \bm{\theta}\in\Gamma,\\
0, & \bm{\theta}\notin\Gamma.
\end{cases}
\]

The truncated prior preserves the relative prior probabilities within $\Gamma$ and assigns zero probability outside $\Gamma$. It does not alter the shape of the prior inside the retained region.

For a single parameter with a unimodal posterior, the excluded values are typically the outer portions of the prior range where the estimated posterior density is very small. These regions are called the tails of the posterior distribution. In multiple dimensions, TMNRE commonly constructs a retained region using bounds obtained from several marginal posterior estimates.

The next round of simulations is drawn from the truncated prior rather than the original prior. Consequently, simulations are concentrated in parameter regions that are more compatible with the observed data. The ratio estimators are retrained using the new simulations, and the truncation can be repeated over several rounds.

The truncation must be conservative. If a genuine posterior region or mode is excluded during an early round, later rounds cannot recover it because no further simulations are drawn there. Thresholds and truncation criteria must therefore be chosen so that regions with non-negligible posterior probability are retained.

TMNRE combines marginal ratio estimation with this iterative restriction of the simulation domain. Marginalisation avoids directly representing an unnecessarily high-dimensional joint posterior, while truncation reduces the number of simulations spent in regions that are unlikely under the observation. The result is a simulation-based inference procedure focused on the marginal posterior quantities required for the analysis.

\end{document}